\section{The UC Proof}\label{app:uc}

This \thispart{} gives the simulator and the hybrid argument behind
Theorem~\ref{thm:uc}. We follow the conventions of the composable ledger
treatment of Badertscher et al.~\cite{bmtz17} and of Generalized
Channels~\cite{genchannels}. $\Ftree$ is protocol-aware: for an honest
\iface{EXIT} or \iface{EVICT} it constructs the transaction of
Figure~\ref{fig:exit} itself, from the honest party's leaf and path and
the fee coin $f$ named in the call, obtains the signature through
$\Fsig$, and submits it to $\Gledger$. The simulator never builds an
honest transaction; it only schedules. The two worlds therefore agree on
the ledger contents and differ only in how the tree state is kept.

\subsection{The Hybrid World}

The real protocol $\Pi_\tau$ runs in the presence of:

\begin{description}
\item[{$\Gledger[\Pred_\tau]$}] the global ledger functionality of
  Badertscher et al., parameterised by a validity predicate that, for a
  transaction spending a tree output, evaluates the tier's covenant
  predicate $\mathsf{Pred}_\tau$ on the input state, the witness and the
  outputs (\S\ref{sec:protocol}). It guarantees persistence, and liveness
  with parameter $\Delta_L$: a valid transaction with sufficient fee that
  an honest party submits, and that stays valid, is in the ledger state of
  every honest party within $\Delta_L$ rounds. The adversary orders
  transactions within the rules, and may get a conflicting spend in first,
  which invalidates the honest one; [C] bounds how often.
\item[$\Gclock$] the global clock of Katz et al.~\cite{kmtz13}; rounds are
  blocks.
\item[$\Fsig$] the ideal signature functionality~\cite{canettiSig}, which
  we use in place of an EUF-CMA reduction; every \texttt{CHECKSIG} and
  \texttt{CHECKSIGFROMSTACK} in the scripts is a call to its verification
  interface, with one exception: MIN's \op{CHECKSIG} under the generator
  $G$, whose key has no honest signer, is real BIP-340 verification, and
  its binding is assumption [MIN] (Table~\ref{tab:assume}).
\end{description}

Hashes are treated in the idealised model of \S\ref{sec:security}
(injectivity plus tag separation), and the taproot commitment in the
random-oracle model where the two internal keys share a sign
(Appendix~\ref{app:lean}). Every abort event below is one in which the
simulator holds two distinct preimages of one digest, so the argument
transfers to a collision-resistance reduction.

\subsection{The Simulator and the Invariant}

$\Sim$ runs an internal copy of $\Adv$ and of every honest party's
protocol state, and keeps a partial tree $\hat T$ equal to $\Ftree$'s.
Claim~\ref{clm:inv} is the invariant: after each ledger event,
$\hat T$ has the root the tree output commits to, and every honest
party's leaf is a known node of $\hat T$ whose path folds to that root.

\paragraph{Ledger events.} Whenever a transaction $\mathsf{tx}$ spending a
tree output appears in the ledger, $\Sim$ parses its witness:

\begin{enumerate}
\item Identify the script leaf from the control block; this fixes the
  transition type and its parameters ($K$, $h$, profile, variant).
\item For each step, fold the supplied node up the supplied path. Where
  the path passes through an opaque node of $\hat T$, call
  $\iface{REVEAL}$ for each level with the nodes the witness supplies
  (they hash to the opaque node because the fold reaches the committed
  root); where it passes through a known node and the fold disagrees with
  it, $\Sim$ holds two preimages of one digest (abort event
  $\mathsf{Coll}$). Otherwise Theorem~\ref{thm:binding} gives the unique
  position $p$ the step proves.
\item Map the steps to the interface: a leaf replaced by $\Emp_0$ is
  \iface{EXIT} (owner flag) or \iface{EVICT} (evictor flag); $\Emp_h$
  replaced by a node is \iface{JOIN}, with the subtree $\Sim$ obtains from
  the operator's join messages, or its opaque root; a leaf replaced by a
  leaf is \iface{UPDATE}; the two-step carve is \iface{VOUCH} followed by
  \iface{CLAIM} (for \tr{SPLIT}, or a voucher a corrupt owner issued) or
  \iface{CLAIM} of a recorded voucher; the post-order pattern is
  \iface{SWEEP}.
\item Check the authorising signatures with $\Fsig$. A valid signature
  under an honest key that the honest party did not produce is a forgery
  (abort event $\mathsf{Forge}$).
\item If $\mathsf{tx}$ is an honest call that $\Ftree$ posted, let it take
  effect; otherwise invoke the interface on behalf of the corrupt party
  that authorised it.
\end{enumerate}

\paragraph{Honest calls.} When an honest party calls an interface,
$\Ftree$ builds and posts the transaction and $\Sim$ chooses, within
$\lat(q)$, when it takes effect, matching the round in which $\Gledger$
includes it. When a corrupt operator grafts or creates a node whose
preimage $\Sim$ cannot obtain, $\hat T$ holds it as opaque; an honest
member never funds such a node without verifying its own path, and
$\Sim$ reveals that path from the member's state, so the member's leaf
is known.

\paragraph{The invariant, by case.} Assume Claim~\ref{clm:inv} before an
accepted spend. \emph{Exit, evict, update, carve.} Each step is valid from
the committed root, so by Lemma~\ref{lem:step} the new root is the root
of $\hat T$ with the step's position replaced, which is $\Ftree$'s tree
after the call; Theorem~\ref{thm:cons} gives that the amounts agree, and
the bytecode that the destinations do. A drained carve source becomes
$\Emp_0$ in both (\lean{claimedSource}). \emph{Join and create.} The new
node is placed at an empty position (\lean{join\_sound}), known or
opaque. \emph{Sweep.} \lean{evictSubtree\_sound} gives the emptied
subtree and one payout per leaf. \emph{Reveal.} It changes no root. In
every case an honest party's leaf outside the changed position keeps a
valid path by Lemma~\ref{lem:path}, and one inside it was removed by its
owner's or evictor's own call.

\paragraph{Proof of Lemma~\ref{lem:path}.} Let the step at $p'$ have old
node $n$, path $\pi_s$ and new node $n'$, and let $k$ be the height at
which the paths of $p$ and $p'$ meet. Both $\fold(\ell, \pi)$ and
$\fold(n, \pi_s)$ equal $r$; reading the two folds from the root down,
injectivity of $\Hnode$ equates, level by level, the two parents'
preimages, so the siblings above height $k$ coincide, $\ell$'s sibling at
height $k - 1$ is the node the step's fold reaches at that height, and
vice versa. Folding $n'$ up $\pi_s$ changes only that node and the nodes
above it; replacing $\ell$'s sibling at that level by the new node and
folding with the unchanged siblings above therefore yields $r'$. Below
height $k - 1$ nothing on $\ell$'s path changed. No real tree is needed,
only the two folds.

\subsection{Hybrids}

\begin{description}
\item[$\mathcal{H}_0$] the real execution of $\Pi_\tau$ with $\Adv$ in
  the $(\Gledger, \Gclock, \Fsig)$-hybrid world.
\item[$\mathcal{H}_1$] as $\mathcal{H}_0$, but the execution aborts on
  $\mathsf{Forge}$. Indistinguishable by the definition of $\Fsig$.
\item[$\mathcal{H}_2$] as $\mathcal{H}_1$, but the execution aborts on
  $\mathsf{Coll}$, on a taproot-commitment collision, and in MIN on a
  violation of [MIN]. In the idealised hash model and the ROM these have
  negligible probability; in each, the simulator outputs the colliding
  preimages.
\item[$\mathcal{H}_3$] as $\mathcal{H}_2$, but the honest parties' tree
  state is replaced by $\hat T$ and their transactions are built and
  posted by $\Ftree$. The invariant above makes the two executions
  identical transaction by transaction; payability
  (Theorem~\ref{thm:pay}, and Lemma~\ref{lem:path} for a path through an
  opaque node) gives that every honest exit $\Ftree$ builds is valid.
\item[$\mathcal{H}_4$] the ideal execution of $\Ftree$ with $\Sim$.
  Identical to $\mathcal{H}_3$ by construction.
\end{description}

\paragraph{Liveness.} $\Ftree$ lets $\Sim$ delay an honest call by up to
$\lat(q)$. In the real world the honest transaction confirms unless a
conflicting spend confirms first; under [C] at most $\kappa$ spends of
the output confirm per block, at most $\beta$ of the conflicting ones are
adversarial, and $q$ honest calls are queued ahead, so after at most
$\lceil (\beta + q)/\kappa \rceil$ blocks of $\lat_L$ rounds each the
honest party's rebuilt transaction, valid by Lemma~\ref{lem:path}, is
included within one further $\lat_L$. Nothing beyond [C] is claimed: the
model does not bound $\beta$, and pinning (\S\ref{sec:security}) is one
way it grows.

\subsection{Per-Tier Remarks}

The argument is the same for every tier because the predicates are the
same; what differs is $\Ftree|_\tau$ (\S\ref{sec:security}) and one
detail each of MIN and CCV. In MIN, \textsf{CHECK\_OLD} and
\textsf{FINAL} are enforced through the constructive signature check of
\S\ref{sec:min}, so $\Sim$ recovers the enforced fields from the witness's
raw pieces rather than from \op{TX} selectors, and their binding is
[MIN]. In CCV the witness carries neither $r$ nor $m$: the input's
commitment to the old root and output~0's to the new one are checked by
\op{CHECKCONTRACTVERIFY} under the input's own taptree, so $\Sim$ reads
$r$ and $m$ from the spent output, and the binding of
\lean{witness\_root\_binds} rests on the BIP-443 data tweak in place of
$\TreeSpk$ (\lean{CCVTweak.binding}, \lean{ccv\_input\_binds}). In every
tier the voucher message names the lineage's $\mathsf{tree\_id}$
(spec~0.10.0), so a voucher for one instance of $\Ftree$ does not verify
in another, and $\Ftree$'s \iface{VOUCH} records it per instance.
\tr{REWRITE}$_h$ and \tr{DETACH}$_h$ are outside the theorem
(\S\ref{sec:security}).

\subsection{Composition for Channel Leaves}

Conjecture~\ref{thm:chan} would follow the pattern of Generalized
Channels~\cite{genchannels}, in the hybrid model with $\Ftree$,
$\Gledger$, $\Gclock$ and $\Fsig$. The channel simulator would reuse the
channel-level simulator of that work for the update and settlement
transactions, which are ordinary ledger transactions once the funding
output exists, and use $\Ftree$'s \iface{EXIT} and \iface{EVICT} to
materialise it. The three tree-layer preconditions of
\S\ref{sec:channels} are the conditions under which $\Ftree$ guarantees
that the materialised output carries the full capacity to the channel
key: exact payouts (\lean{ExitK.exact\_conservation}), non-claimability
(the \lean{SEQ\_FINAL} facts), and a lease at or below the tip so that
\iface{EVICT} is available to the counterparty. The missing pieces are
the HTLC deadline argument under [C] and the rebinding of pre-signed
states across tree transitions; they are exercised block by block on the
simulator's chain model (\lifted{app:channels}), not proved.
